\section{模仿学习的基本思想}

模仿学习（Imitation Learning）起点是直接复制专家，而非设计复杂的奖励函数。它适合数据充足、专家行为质量高但 reward 不明确的场景。

\begin{importantbox}{模仿学习的目标}
给定专家策略 $\pi_{\text{expert}}$ 产生的演示轨迹集 $\mathcal{D} = \{(s_1,a_1,\ldots,s_T)\}$，目标是找到参数化策略 $\pi_\theta$，使其预测的动作分布在相同状态下与专家尽量重合，从而保持专家级性能。
\end{importantbox}

\subsection{Version 0：确定性策略回归}

最直接的做法是把模仿学习当成行为映射：
$$\min_\theta \frac{1}{|\mathcal{D}|}\sum_{(s,a) \in \mathcal{D}} \|\pi_\theta(s) - a\|^2.$$
这种方法完全依赖平衡好的演示数据，对 noise 非常敏感。

\begin{warningbox}{均值问题与不可逆行为}
当专家数据呈现多模态（如左转/右转）时，回归模型会输出“平均”动作，造成在环境中进入危险状态；此问题在城市驾驶、双臂协作等任务中普遍存在。
\end{warningbox}

\subsection{行为克隆 vs. 逆向强化学习}

行为克隆直接建模动作，而逆向强化学习（Inverse RL）则重构 reward，再做 RL 优化。下面对比两个方向的 trade-off：

\begin{table}[H]
\centering
\begin{tabular}{p{0.25\textwidth} p{0.32\textwidth} p{0.32\textwidth}}
\toprule
方法 & 优点 & 缺点 \\
\midrule
行为克隆 & 训练简单，直接使用 supervised learning & 容易失去多模态、受分布偏移影响大 \\
逆向 RL & 能恢复潜在 reward，具备泛化能力 & 需要大量采样，reward 学习与 policy 训练耦合 \\
\bottomrule
\end{tabular}
\caption{行为克隆与逆向 RL 的比较}
\end{table}

\begin{knowledgebox}{为何优先行为克隆}
在很多工业场景（如自动驾驶）中，专家轨迹足够多且高质量，此时直接做行为克隆可以极大降低上线周期；如果出现分布漂移，再逐步引入 reward 模型。
\end{knowledgebox}

\subsection{真实场景中的落地与 caution}

实际部署时常见的策略包括：行驶数据的上采样、以 experts 为 baseline 训练 fallback policy，以及建立 human-in-the-loop 的监控机制。在低频故障场景下，先用行为克隆做 warm-start，再用 RL 微调是普遍设计。

\begin{importantbox}{行为克隆的可解释性}
因为输出是 deterministic/log-prob，方便进行 trace-back；一旦模型输出异常，工程团队可以直接查看最近的 expert sample，快速定位问题源。
\end{importantbox}

%-\subsection{本章小结}
-
%-行为克隆提供了...
%-\subsection{多专家融合与动态权重}
-
-在大模型时代...

\begin{knowledgebox}{专家混合的价值}
将多个专家的 logits 放入 Mixture-of-Experts 框架，可以在少量样本上实现 zero-shot 迁移，因为 gating network 可以自动偏向与当前任务更相关的 expert。
\end{knowledgebox}

\subsection{本章小结}

行为克隆提供了最快速的入门路径、但 deterministic 回归必须配合多专家 gating 与分布建模才能在复杂场景中泛化。多专家混合提供了针对不同 skill 的自动调度，延展了模仿学习的适用边界。

